\section{Health and safety}\label{sec:safety}

Hands-on work starts only after the departmental inductions and with signed risk assessments for the
laser work, the gas set-up, and the chemicals (COSHH). The main hazards and controls are:

\begin{description}[style=nextline,leftmargin=1.2em,itemsep=3pt]
  \item[Near-UV laser radiation] A \SI{375}{nm} diode of tens of milliwatts, or the \SI{369.5}{nm}
        systems, are Class~3B. Near-UV beams are barely visible, so stray reflections are easy to
        miss. Controls: enclose the beam paths of the testbed (the long runs are unattended), fit an
        interlock to the enclosure, wear eyewear with adequate optical density at the wavelength in
        use, use fluorescent viewing cards, register the set-up with the laser safety officer and follow
        the local rules. Unattended operation needs its own risk assessment.
  \item[Compressed and inert gases] Cylinders are secured and fitted with suitable regulators and
        relief devices. The lab \Ntwo{} supply is used where possible. The purge flows are tiny
        (about \SI{0.1}{L/min} in total) and cannot deplete the oxygen in a ventilated lab. The relevant
        scenario is the release of a whole cylinder in a small room. The risk assessment must check the
        room volume and ventilation and decide whether an oxygen-depletion monitor is needed. Exhausts are
        routed to extraction.
  \item[Forming gas] \SI{5}{\percent} \Htwo{} in \Ntwo{} is normally supplied and classified as a
        non-flammable mixture. It is still handled as a hydrogen-containing gas (no ignition sources at
        the exhaust, leak checks with a hydrogen detector). Use only if Campaign~II needs it.
  \item[Oxygen absorbers] Iron-based sachets heat up as they react. Keep them sealed until use, do not
        cut them open, and allow used sachets to finish reacting before disposal.
  \item[Chemicals] Isopropanol and acetone (flammable), epoxies and resins (skin sensitisers),
        polishing slurries. Wear gloves and eye protection and work with ventilation as stated in the
        COSHH assessments.
  \item[3D printing] Run printers in ventilated or filtered enclosures (FDM printers emit ultrafine
        particles and VOCs). Handle uncured resin with gloves.
  \item[Fibre shards] Wear eye protection when cleaving and polishing, use a sharps container and
        a dedicated work mat, and keep food and drink out of the area.
  \item[Plasma and UV/ozone cleaners] If used, they produce UV-C and ozone. Use them only with their
        interlocks and extraction.
\end{description}

\begin{safetybox}[Rule of thumb]
If an experiment is going to run unattended for weeks, design it so that every failure (power cut, gas
cylinder empty, laser fault, interlock trip) leaves it in a safe state \emph{and} is logged with a
timestamp. Then the failure costs you a data gap but does not put anyone at risk.
\end{safetybox}
